%!TEX root = ../template.tex
%%%%%%%%%%%%%%%%%%%%%%%%%%%%%%%%%%%%%%%%%%%%%%%%%%%%%%%%%%%%%%%%%%%%
%% 5-Conlcusions.tex
%% NOVA thesis document file
%%
%% Chapter with the template manual
%%%%%%%%%%%%%%%%%%%%%%%%%%%%%%%%%%%%%%%%%%%%%%%%%%%%%%%%%%%%%%%%%%%%

\typeout{NT FILE 7-Conlcusions.tex}%

\chapter{Conclusions and Future Work}
\label{cha:Conclusions and Future Work}

\glsresetall

\section{Conclusion}

This thesis demonstrates, both in simulation and experimentally, the feasibility of using drones to encircle a target in order to discourage its entrance into a restricted aerial space.

\section{Future Work}

Although the objectives of this thesis were achieved, several directions for future research remain open. One possible extension is the inclusion of physical obstacles in the environment, such as buildings or trees, while guaranteeing collision avoidance with both static and dynamic obstacles. This could be achieved through the integration of LIDAR measurements or other sensing and perception methods.

Another important improvement would be the use of more realistic target localization methods, such as radar-based position measurements, instead of idealized sensing assumptions.

It may also be worthwhile to investigate alternative agent distributions within the mesh beyond the uniformly spaced configuration considered in this work. If the most likely direction of motion of the target is known, it may be advantageous to concentrate more agents in that direction or to deform the mesh accordingly. One possible approach is the use of weighted Voronoi diagrams.

Finally, future work could explore making the encirclement distance between the target and the swarm dependent on the reachability sets of either the target or the swarm agents. For example, given the dynamic characteristics of the drones and a desired overlap between the agents' reachability sets (e.g., 10\%), it would be useful to determine the maximum allowable spacing between agents and the target while still guaranteeing the required overlap.

Professor eu sei que está curto mas é para ter algo escrito, ainda não acabei a tese por isso vou deixar isto muito breve até Agosto/Setembro.
